\paragraph{Local estimates are not distributed admission control.}
An independent frontend cannot see another frontend's reservations before the
next observation. Several frontends can therefore choose the same apparently
idle replica, and a domain can grow later than the actual arrival burst.
Deterministic primary maps reduce unnecessary replication but can also correlate
those decisions. This design offers no global queue bound, no per-adapter
starvation guarantee, and no theorem of throughput optimality. Obtaining such
properties would require an admission protocol, migration, or explicit engine
scheduling support beyond this prototype.

The preferred map also counts registered versions that currently have no demand.
Long-running training can accumulate a large inactive catalog even when only a
few latest versions are sampled. The evaluated catalog contains at most four
versions per job and 64 versions overall. Demand-filtered catalog retention or
version garbage collection would require an additional policy; the present
results do not establish behavior under indefinite publication history.

The cache exception has another explicit limit: if every replica advertises
some reusable prefix for an adapter, every replica passes the locality filter.
The work score still compares the remaining prefill, but domain restriction no
longer limits that adapter's replication. A common system prefix across many
otherwise different prompts could produce this condition. Our controlled
problem prefixes are distinct, so the experiments do not resolve this case.

For example, consider eight frontends with identical empty observations, each
submitting 36 requests for the same uncached adapter before a refresh. Each
frontend sees one batch of demand and selects the same primary. Executing this
case against the actual router on CPU assigns all 288 requests to one replica.
This is a deterministic counterexample, not a GPU performance measurement. The
local score bound does not prevent this correlated-arrival case. A batched,
replica-local admission protocol could address it without a cluster-wide request
coordinator, but it is not implemented or evaluated here.

\paragraph{Autoscaling remains a separate control problem.}
Flash scales on request concurrency, whereas the router compares estimated work
and locality. A request-count target and an estimated-work budget can disagree,
especially with heterogeneous output lengths. The fixed-size measurements can
identify idle replicas during pressure, but cannot establish that scale-out
occurs only after all existing replicas are saturated. Membership retention
reduces remapping after a missed poll; it does not validate scaling stability.
An elasticity experiment with changing replica count is necessary before using
these results to tune Flash's autoscaling thresholds.

\paragraph{Workload and model scope.}
The experiments cover one hybrid 27B model, one adapter rank, one GPU type, and
eight replicas. Long inputs and exact decode lengths make comparisons
interpretable, but omit natural language structure, early EOS, variable rank,
and long-tailed output distributions. Saved checkpoint arrivals are real adapter
registrations, while their optimizer delays are replayed. We do not infer
training convergence, reward quality, or optimizer throughput from serving
measurements. Larger-model tensor parallelism and heterogeneous devices remain
untested.
The eight-slot GPU adapter budget is also fixed. Increasing it could remove
some of the measured adapter interference at the cost of memory available for
other state. We do not evaluate that cache-capacity tradeoff or claim the same
gain when every replica can retain the full adapter catalog.

\paragraph{What a cache hit does and does not show.}
A high requested-token cache ratio after auxiliary preparation is expected.
The seed still performed the cold computation. Similarly, more first-prefill
hot-adapter hits need not translate to fewer later GPU loads, and summed
host-side loading time is not a wall-clock latency decomposition. This is why
the evaluation reports auxiliary uncached tokens, actual load counters,
completion latency, and useful output throughput together. An optimization is
not successful merely because one locality metric improves.

\paragraph{Transport, observation, and engine interactions.}
An HTTP timeout can occur after engine admission. Retrying then creates work
that is absent from the logical request count. Our stress pilots retain that
cost. Increasing the timeout did not eliminate a separate HTTP/1 read-failure
mode; the subsequent controlled comparisons use HTTP/2 for every policy.
Sidecar presence alone cannot distinguish a request waiting for an
engine connection from one queued inside the engine. Longest-prefix-match
scheduling, the sidecar's connection pool, cache eviction, and Flash transport
all affect end-to-end behavior. The available instrumentation does not provide a
complete causal decomposition of every tail event.

\paragraph{Evidence strength and deployment status.}
Pilot-driven design choices are evaluated separately from the later fixed
workloads. Repeated trials are reported individually; two executions cannot
establish a population confidence interval or robustness to all production
traffic. Literature comparisons describe mechanisms rather than head-to-head
reimplementations. Domains is available through the explicit
\texttt{inference\_lora\_routing} recipe option, while cohort preparation remains
a benchmark helper. The production default is unchanged. These results support
the measured conclusions, not a general claim of state-of-the-art serving
performance.
